% BEGIN REV09_PELL_CATALAN
\subsection{Pell dilation, Catalan moment and oriented return}
\label{corpus9:iv:pell}

The preceding parametric moment \(\mathcal C_2\) retains the
depth variable; write
\(G_{\mathrm{Cat}}=\mathcal C_2(1)\).
Its following specialization receives
the scale of the Pell spectral screw
\[
 V^{-1}\mathscr S V=\lambda_P i\,\mathscr S,\qquad
 \lambda_P=2+\sqrt3,\qquad \rho=\lambda_P^{-1}=2-\sqrt3.
\]
The covariance law is the operator hypothesis of this
specialization; it is understood on the common invariant
domain of \(\mathscr S\) and the invertible transport \(V\).
The unit \(\lambda_P\) belongs to \(\mathbb Q(\sqrt3)\);
the unit \(3+2\sqrt2\) of the compression realization
belongs to \(\mathbb Q(\sqrt2)\) and retains its own reader.
On its scalar transverse orbit,
\[
 z_n=z_0(\lambda_P i)^n,\qquad
 r_n=r_0\lambda_P^n,\qquad
 \theta_n=\theta_0+\frac{\pi_{\mathrm{HMT}}n}{2}.
\]
For \(z_0\ne0\), eliminating \(n\) gives the suspension
\[
 r(\theta)=r_0\exp\!\left(
   \frac{2\log\lambda_P}{\pi_{\mathrm{HMT}}}
   (\theta-\theta_0)\right).
\]
The inverse branch has multiplier \(-i\rho\). This
scale realization receives the screw as its operator
antecedent and retains its orientation.

\begin{proposition}[Spectral reading and evaluation at the contractive scale]
\label{corpus9:iv:pell-evaluacion}
On \(\mathcal H=\ell^2(\mathbb N_{>0})\), let
\[
 N\delta_n=n\delta_n,\qquad U_4\delta_n=i^n\delta_n,\qquad
 Q_4=\frac{U_4-U_4^*}{2i},\quad
 C_4=\frac{U_4+U_4^*}{2}.
\]
For \(0<r<1\) and \(\nu\in\mathbb C\), the traces
\[
 F_\nu(r)=\operatorname{Tr}(N^{-\nu}r^NQ_4),\qquad
 H_\nu(r)=\operatorname{Tr}(N^{-\nu}r^NC_4)
\]
converge absolutely and locally uniformly in
\((\nu,r)\). One has \(F_2(r)=\mathcal C_2(r)\) and
\begin{align}
 \mathcal C_2(\rho)
   &=\frac23G_{\mathrm{Cat}}
     -\frac{\pi_{\mathrm{HMT}}}{12}\log\lambda_P,\\
 \mathcal D\mathcal C_2(\rho)&=\frac{\pi_{\mathrm{HMT}}}{12},
 &\mathcal D^2\mathcal C_2(\rho)&=\frac14 .
 \label{corpus9:iv:pell-valores}
\end{align}
\end{proposition}
\begin{proof}
The diagonal basis gives the series with coefficients
\(\sin(n\pi/2)\) and \(\cos(n\pi/2)\). On a compact set of
\(\nu\) and with \(r\leq r_0<1\), their absolute values are
bounded by \(n^M r_0^n\) for some \(M\), a summable series.
Selection of the odd indices in the first trace
produces exactly \(\mathcal C_2(r)\).

The tangent subtraction formula gives
\(\rho=\tan(\pi_{\mathrm{HMT}}/12)\), as a subsequent
angular recognition of the Pell unit. For
\(0<\theta<\pi/2\), integration by parts of the moment gives
\[
 \mathcal C_2(\tan\theta)
  =\theta\log(\tan\theta)-\int_0^\theta\log(\tan x)\,dx.
\]
The boundary term at zero vanishes. Subtracting the
Fourier series of \(\log(2\sin x)\) and \(\log(2\cos x)\),
and integrating with Abel regularization, gives
\[
 -\int_0^\theta\log(\tan x)\,dx
   =\sum_{\substack{n\geq1\\2\nmid n}}
          \frac{\sin(2n\theta)}{n^2}.
\]
The regularization converges in the integral because the
logarithmic singularity at zero is integrable; the
integrated series is dominated by \(\sum n^{-2}\).

For odd \(n\), let \(\chi_{-4}(n)=(-1)^{(n-1)/2}\).
The identity, checked in the six odd residue classes modulo
twelve, is
\[
 2\sin(n\pi_{\mathrm{HMT}}/6)
   =\chi_{-4}(n)
     +3\,\mathbf1_{3\mid n}\chi_{-4}(n/3).
\]
Its sum with weight \(n^{-2}\) gives
\(\frac12(1+3/9)G_{\mathrm{Cat}}=2G_{\mathrm{Cat}}/3\).
Taking \(\theta=\pi_{\mathrm{HMT}}/12\) and
\(\log\rho=-\log\lambda_P\) yields the first
evaluation. The others follow from
\(\mathcal D\mathcal C_2(r)=\arctan r\),
\(\mathcal D^2\mathcal C_2(r)=r/(1+r^2)\) and
\(\rho+\rho^{-1}=4\).
\end{proof}

The real component of the same trace retains
\[
 H_2(r)=\frac14\sum_{k\geq1}\frac{(-r^2)^k}{k^2},
 \qquad
 \mathcal D H_2(r)=-\frac12\log(1+r^2),\qquad
 \mathcal D^2H_2(r)=-\frac{r^2}{1+r^2}.
\]
These formulas follow by selecting the even indices and
differentiating the absolutely convergent series. Hence
\[
 \operatorname{Tr}(N^{-2}r^NU_4)=H_2(r)+i\mathcal C_2(r).
\]
The designation
\(\operatorname{Li}_2(z)=\sum_{n\geq1}z^n/n^2\)
subsequently recognizes this trace as
\(\operatorname{Li}_2(ir)\). In particular,
\[
 \operatorname{Tr}(N^{-2}\rho^NU_4)
 =\frac14\operatorname{Li}_2(-\rho^2)
  +i\left(\frac23G_{\mathrm{Cat}}
       -\frac{\pi_{\mathrm{HMT}}}{12}\log\lambda_P\right).
\]
Replacing \(U_4\) by \(U_4^*\) preserves the real
part and reverses the imaginary part.

\begin{proposition}[Composition of depth and orientation]
\label{corpus9:iv:pell-composicion}
For \(0<r,s<1\), \(a,b\in\mathbb Z/4\mathbb Z\) and \(m\geq0\), let
\(B(r,a)=r^NU_4^a\). Then
\[
 B(r,a)B(s,b)=B(rs,a+b),\qquad
 B(r,1)^m=r^{mN}U_4^m.
\]
In particular, \(B(r,1)^4=r^{4N}\ne I\) for \(0<r<1\).
\end{proposition}
\begin{proof}
Both factors are diagonal in the same basis. On
\(\delta_n\), their product multiplies by
\(r^ni^{an}s^ni^{bn}=(rs)^ni^{(a+b)n}\).
The power formula follows by induction, and
\(U_4^4=I\) gives the fourth return. The factor
\(r^{4N}\) remains distinct from the identity.
\end{proof}

The power reader of the inverse multiplier
\(-i\rho\) is \(B(\rho,-1)\). Its imaginary moment has
the opposite orientation to that of \(B(\rho,1)\), whereas
the damped depth is the same. The transport \(V\),
observable \(\mathscr S\) and reader \(U_4\) retain
their respective spaces: the preceding composition concerns
the orbit multipliers and their diagonal reading.
Phase returns after four steps, while accumulated
contraction remains recorded.
% END REV09_PELL_CATALAN
